% Generated draft: requires mathematical and source review.
\documentclass[11pt]{article}
\usepackage[margin=25mm]{geometry}
\usepackage{fontspec}
\setmainfont{Latin Modern Roman}
\usepackage{amsmath,amssymb,mathtools,unicode-math}
\setmathfont{Latin Modern Math}
\usepackage{enumitem,needspace,xurl,hyperref}
\setlength{\parindent}{0pt}
\setlength{\parskip}{5pt}
\setlength{\emergencystretch}{3em}
\newcommand{\N}{\mathbb{N}}
\newcommand{\Z}{\mathbb{Z}}
\newcommand{\Q}{\mathbb{Q}}
\newcommand{\R}{\mathbb{R}}
\newcommand{\C}{\mathbb{C}}
\newcommand{\F}{\mathbb{F}}
\newcommand{\A}{\mathbb{A}}
\newcommand{\PP}{\mathbb{P}}
\newcommand{\E}{\mathbb{E}}
\newcommand{\eps}{\varepsilon}
\newcommand{\dd}{\,\mathrm{d}}
\begin{document}
\section*{UnsolvedMath ID: 10400062}
\textbf{Source problem code:} catalogue entry
\subsection{Definitions and mathematical statement}
A zero framing of an oriented knot \(K\subset S^3\) is a framing for which a push-off has linking number zero with \(K\). For a geometric representative \(\gamma\) in \(\R^3\), the self Gauss integral is conventionally the writhe integral
\[
 \operatorname{Gauss}(\gamma)=\frac1{4\pi}
 \int_{S^1\times S^1\setminus\Delta}
 \frac{(\gamma'(s)\times\gamma'(t))\cdot(\gamma(s)-\gamma(t))}
 {\|\gamma(s)-\gamma(t)\|^3}\,\dd s\,\dd t.
\]
Zero framing and vanishing of this geometric integral are distinct conditions and are both retained.

C. Lescop asks whether, for the same zero-framed knot represented with \(\operatorname{Gauss}(\gamma)=0\), the Kontsevich integral equals the Chern--Simons series obtained from configuration-space integrals:
\[
 Z_{\mathrm{Kont}}(K)=Z_{\mathrm{CS}}(\gamma).
\]
The intended equality is of the full diagram-valued formal series in a common degree-completed diagram space, hence coefficientwise in every degree. Equality only after a chosen Lie-algebra weight system or an unspecified reparametrization would be a different assertion.

The excerpt does not define the particular configuration-space series, its anomaly or correction conventions, the diagram coefficients, or normalization of either side. These must be recovered from the source; they cannot be freely adjusted to force equality. The Gauss-integral display gives its conventional geometric meaning, which also needs confirmation against the source conventions.
\subsection{Short English statement}
For a zero-framed knot with vanishing Gauss integral, is its Kontsevich integral exactly the Chern--Simons configuration-space integral series? The precise series and normalizations need source definitions, and zero framing must not be conflated with zero writhe.
\subsection{Sources}
\begin{itemize}
\item Local UnsolvedMath registry, permanent ID 10400062, source code catalogue entry; statement attribution C. Lescop; proposer field and proposal year absent; update 2026-08-06.
\item Ohtsuki, \emph{Problems on invariants of knots and 3-manifolds} (2002), Question 3.12, PDF page 57: \url{https://msp.org/gtm/2002/04/gtm-2002-04-024s.pdf}. Registry access date 2026-07-29; PDF font-structure extraction. Provenance distinguishes 231 numbered bold headings. Present primary-literature status and redistribution terms require review.
\item The supplied summary repeats the comparison problem without an exact configuration-integral reference or verified resolution.
\end{itemize}
\end{document}
